\chapter{Konstruktionsbegründung}
\label{chap:konstruktionsbegruendung}

\section{Positionierung von Radiator und Lüfter}
\label{sec:Positionierung}

Für die Anordnung der Radiatoren am Fahrzeug existieren mehrere Möglichkeiten,
die jeweils unterschiedliche Kompromisse erfordern. Die Wahl der Position
beeinflusst die Qualität der Anströmung, die Aerodynamik des Gesamtfahrzeugs,
die Lage des Schwerpunkts sowie den Konstruktions- und Fertigungsaufwand.

\subsection{Randbedingungen}
\label{subsec:randbedingungen_position}

Folgende Randbedingungen waren bei der Positionierung zu berücksichtigen:

\begin{itemize}
    \item \textbf{Anströmung:} Je weiter vorne der Radiator sitzt, desto
          sauberer ist die anströmende Luft. Gleichzeitig verwirbelt er die
          Luft für alle dahinterliegenden Bauteile stärker, insbesondere für
          die Aerodynamik im hinteren Fahrzeugbereich.
    \item \textbf{Einbauhöhe:} Höhere Positionen ermöglichen in der Regel
          eine sauberere Anströmung, tiefere Positionen sind hingegen
          deutlich günstiger für die Schwerpunkthöhe.
    \item \textbf{Fahrersichtfeld und Driver Fluid Protection:} Nach
          Regel T0.0.00 dürfen im Sichtfeld des Fahrers keine
          wasserführenden Schläuche verlaufen. Zusätzlich ist die Driver
          Fluid Protection (T0.0.0) einzuhalten.
    \item \textbf{Trägheit:} Massen, die weit von der Hochachse entfernt
          montiert werden, erhöhen das Trägheitsmoment des Fahrzeugs.
\end{itemize}

\subsection{Betrachtete Positionsvarianten}
\label{subsec:positionsvarianten}

In den letzten Jahren haben sich bei den Teams drei grundsätzliche Positionen
durchgesetzt, die jeweils mehrere Untervarianten aufweisen.

\paragraph{Position im Frontwing.}
Diese Variante wird beispielsweise vom Team Amberg-Weiden Running Snails
eingesetzt~\cite{TODO_Quelle_Running_Snails}. Die Radiatoren sind dabei in
einem neu gestalteten Frontwing untergebracht, und ein Cooling Duct verläuft
durch den Frontwing. Laut Reglement wird dadurch ein sonst ungenutzter
Bauraum verwendet, und der Radiator erhält die beste und sauberste
Anströmung. Nachteilig sind die relativ hohe Schwerpunktlage sowie ein
erhöhtes Trägheitsmoment des Fahrzeugs, da die Masse weit entfernt vom
Drehpunkt montiert wird. Auch die Einhaltung der Driver Fluid Protection
(T0.0.0) ist bei dieser Variante anspruchsvoll.

\paragraph{Position im Seitenkasten.}
Ein Beispiel für diese Variante ist der RP25e des Dynamics e.V. Die
Radiatoren werden typischerweise relativ weit hinten im Seitenkasten an den
Endplates montiert. Das stellt einen guten Kompromiss dar zwischen der
Frischluft, die den Radiator erreicht, und der Verwirbelung, die andere
Systeme beeinflussen kann. Die Driver Fluid Protection ist mit mittlerem
Aufwand lösbar, und der Abstand zur Hochachse ist vergleichsweise gering.
Herausforderungen liegen in der Anbindung, da die Aerodynamikentwicklung nur
wenig Bauraum zur Verfügung stellt, sowie in der durch das Reglement stark
begrenzten vertikalen Bauhöhe. Für diese Position liegen außerdem
Volumenstromdaten aus den vergangenen Saisons vor.

\paragraph{Position am Fahrzeugheck.}
Bei dieser Variante wird an der Anti Collision Plate (ACP) eine zusätzliche
Haube angebracht, deren alleinige Aufgabe die Halterung der Radiatoren ist.
Hinsichtlich des Bauraums ist dies die freieste Position, da kein anderes
Subsystem sie mitnutzt. Der natürliche Luftstrom durch den Fahrtwind ist
hier jedoch nahezu nicht vorhanden, sodass der Einsatz von Lüftern
unerlässlich ist. Die Driver Fluid Protection lässt sich ohne nennenswerten
Mehraufwand erfüllen, da Seitenkasten und weitere Aerodynamikkomponenten die
Schläuche fast vollständig abdecken. Nachteilig ist der Mehraufwand für
Konstruktion und Fertigung der ein bis zwei Lüfterhauben. Zudem müssen diese
so ausgelegt sein, dass die ACP jederzeit demontiert werden kann, da dies
für den Ein- und Ausbau des Akkus zum Laden des Fahrzeugs erforderlich ist.
Daher ist eine erhöhte Anzahl an Schnellverschlüssen bzw. an schnell
steckbaren Verbindungen notwendig.

\begin{table}[htbp]
    \centering
    \caption{Vergleich der betrachteten Radiatorpositionen}
    \label{tab:positionsvergleich}
    \begin{tabular}{@{}p{0.16\textwidth}p{0.38\textwidth}p{0.38\textwidth}@{}}
        \hline
        \textbf{Position} & \textbf{Vorteile} & \textbf{Nachteile} \\
        \hline
        Frontwing
            & beste und sauberste Anströmung; nutzt sonst ungenutzten Bauraum
            & hoher Schwerpunkt; erhöhtes Trägheitsmoment; Driver Fluid
              Protection anspruchsvoll \\
        \hline
        Seitenkasten
            & guter Kompromiss aus Anströmung und Verwirbelung; geringer
              Abstand zur Hochachse; Driver Fluid Protection mit mittlerem
              Aufwand lösbar; Volumenstromdaten vorhanden
            & wenig Bauraum und begrenzte Bauhöhe; aufwendige Anbindung \\
        \hline
        Heck (ACP)
            & größter Bauraum; Driver Fluid Protection sehr einfach
            & kaum natürlicher Luftstrom (Lüfter zwingend); Mehraufwand für
              Lüfterhauben; ACP muss demontierbar bleiben \\
        \hline
    \end{tabular}
\end{table}

\subsection{Entscheidung}
\label{subsec:entscheidung_position}

Nach Abwägung aller Einflussfaktoren wurde die Position im Seitenkasten
gewählt. Um sicherzustellen, dass der benötigte Bauraum zur Verfügung steht
und die aerodynamischen Nachteile begrenzt bleiben, wurde die Auslegung eng
mit dem Aerodynamikteam abgestimmt. Die grobe Einbaulage ist in
Abbildung~\ref{fig:einbaulage_radiator} dargestellt.

% FIGURE-VORSCHLAG A: CAD-Übersicht des Fahrzeugs mit hervorgehobener
% Radiator-/Lüfterposition im Seitenkasten (Seitenansicht oder Iso-Ansicht).
\begin{figure}[htbp]
    \centering
    % \includegraphics[width=\textwidth]{bilder/einbaulage_radiator}
    \caption{Grobe Einbaulage von Radiator und Lüfter im Seitenkasten}
    \label{fig:einbaulage_radiator}
\end{figure}

\subsection{Umsetzung im Seitenkasten}
\label{subsec:umsetzung_seitenkasten}

\paragraph{Schrägstellung des Radiators.}
Die in Kapitel~\ref{chap:UNBEKANNT} festgelegten Radiatoren sind
breiter als der im Seitenkasten verfügbare Bauraum. Der Radiator wird daher
in einem Winkel montiert. Der Kompromiss wurde hierbei bewusst zugunsten der
Aerodynamik gelegt, sodass die Fertigung eines Sonderflügelprofils nicht
erforderlich ist.

\paragraph{Lüfterhaube.}
Durch die Schrägstellung gestaltet sich auch die Positionierung der Lüfter
schwierig. Deshalb wurde eine einfache Haube konstruiert, die es erlaubt, die
Lüfter in einem gewissen Abstand zum Radiator anzuordnen. Dadurch kann die
Abluft direkt nach außen und vom Reifen weg geführt werden. Zudem soll der
an der Fahrzeugaußenseite entstehende Unterdruck genutzt werden, um die
Durchströmung zu verbessern und den Volumenstrom durch den Radiator zu
erhöhen.

% FIGURE-VORSCHLAG B: CAD-Schnitt oder Draufsicht des Seitenkastens mit
% schräg montiertem Radiator, Lüfterhaube und Pfeilen für Zuluft/Abluft.
\begin{figure}[htbp]
    \centering
    % \includegraphics[width=\textwidth]{bilder/radiator_luefterhaube}
    \caption{Schräg montierter Radiator mit Lüfterhaube im Seitenkasten
             (Strömungsrichtung der Luft eingezeichnet)}
    \label{fig:radiator_luefterhaube}
\end{figure}

\paragraph{Lüfterposition und Finger Guards.}
Die Lüfterpositionen ergeben sich maßgeblich aus der Lage von Radiator und
Haube. Gemäß Regel T0.0.0 dürfen sich keine Finger in rotierende
Teile einführen lassen, sobald der Durchmesser der Öffnungen einen
bestimmten Wert überschreitet. Zur Erfüllung dieser Regel wurden eigene
Finger Guards konstruiert und ausgelegt.

% FIGURE-VORSCHLAG C: Foto oder CAD der selbstgebauten Finger Guards am Lüfter.
\begin{figure}[htbp]
    \centering
    % \includegraphics[width=0.7\textwidth]{bilder/finger_guard}
    \caption{Selbst konstruierter Finger Guard am Lüfter}
    \label{fig:finger_guard}
\end{figure}

\section{Positionierung der Pumpen}
\label{sec:pumpen_position}

Die Pumpen lassen sich vergleichsweise frei positionieren. Um das Packaging
zu vereinfachen, wurde entschieden, sie nicht im Fahrzeuginneren
anzuordnen, sondern in die Seitenkastenprofile zu verlagern. Hierfür wurden
einfache Halter ausgelegt, die in die Flügelprofile eingeklebt werden und
die Pumpen allein mittels Kabelbinder fixieren. Diese Lösung ist
äußerst gewichtssparend, was für diese Komponente aufgrund der
reglementbedingten Vorgaben ein wesentliches Ziel war.

% FIGURE-VORSCHLAG D: Foto/CAD der in das Seitenkastenprofil eingeklebten
% Pumpenhalter mit montierter Pumpe.
\begin{figure}[htbp]
    \centering
    % \includegraphics[width=0.7\textwidth]{bilder/pumpenhalter}
    \caption{Pumpenhalter im Seitenkastenprofil}
    \label{fig:pumpenhalter}
\end{figure}

\section{Sensorhalterungen}
\label{sec:sensorhalterungen}

Für die Sensoren wurden neue Halterungen konstruiert, die sowohl am
Prüfstand als auch am Fahrzeug verwendet werden können. Sie wurden aus
Kostengründen und aus Gründen der Isolation im FDM-Verfahren aus Kunststoff
gedruckt. Eine erste Version aus SLA-Harz erwies sich als nicht wasserfest
und wurde unter Last gummiartig weich. Die aktuellen Teile bestehen aus
einem Spezialfilament (Materialbezeichnung), mit dem diese Probleme
vollständig beseitigt werden konnten.

% FIGURE-VORSCHLAG E: Foto der FDM-Sensorhalter (ggf. neben dem ersten
% SLA-Teil zum Vergleich) sowie im eingebauten Zustand am Kühlkreis.
\begin{figure}[htbp]
    \centering
    % \includegraphics[width=0.7\textwidth]{bilder/sensorhalter}
    \caption{FDM-gedruckte Sensorhalter für Prüfstand und Fahrzeug}
    \label{fig:sensorhalter}
\end{figure}

\section{Entlüftungskonzept}
\label{sec:entlueftung}

Für ein schnelles und zuverlässiges Entlüften des Kühlkreislaufs wurde ein
eigenes Entlüftungskonzept entwickelt. Es besteht aus drei Maßnahmen:
den Entlüftungsadaptern, der geänderten Anordnung von Pumpe und Radiator
sowie einer Entlüftungsvorrichtung.

\subsection{Entlüftungsadapter}
\label{subsec:entlueftungsadapter}

Je Kühlkreis kommt ein Krontec-Adapter~\cite{TODO_Datenblatt_Krontec} zum
Einsatz. Er ermöglicht das Öffnen und Schließen des Kreises, ohne dass Wasser
austritt oder Luft eintritt. Die Adapter sind an den höchsten Stellen der
Kühlkreisläufe positioniert, nämlich direkt nach den Radiatoren. Dort sammelt
sich die eingeschlossene Luft und kann beim Entlüften direkt nach außen
gespült werden.

% FIGURE-VORSCHLAG 16: Foto des Krontec-Adapters im eingebauten Zustand
% oder CAD-Ausschnitt mit Markierung der Einbauposition am höchsten Punkt.

\subsection{Anordnung von Pumpe und Radiator}
\label{subsec:anordnung_pumpe_radiator}

Zusätzlich wurden die Positionen von Pumpe und Radiator im Kreis getauscht.
Dadurch liegt der Entlüftungsanschluss vor der Pumpe und nach dem Radiator.
Es muss nur ein kurzes Schlauchstück vorgefüllt werden, bis die Pumpe
genügend Wasser fördert, um die restliche Luft aus dem System zu pumpen.

% FIGURE-VORSCHLAG 17: Schematischer Kühlkreislauf (Fließbild) vorher/nachher
% mit eingezeichneter Reihenfolge von Radiator, Pumpe und Entlüftungsanschluss.
\begin{figure}[htbp]
    \centering
    % \includegraphics[width=\textwidth]{bilder/kuehlkreis_vorher_nachher}
    \caption{Anordnung von Pumpe, Radiator und Entlüftungsanschluss vor und
             nach der Umstellung}
    \label{fig:kuehlkreis_layout}
\end{figure}

\subsection{Vorbefüllen}
\label{subsec:vorbefuellen}

Zum Vorbefüllen kann entweder ein Trichter angeschlossen werden, oder die
Dosieröffnung des Wasserbehälters wird aufgeschnitten, sodass das System
vorläufig direkt daraus befüllt werden kann.

\subsection{Entlüftungsvorrichtung}
\label{subsec:entlueftungsvorrichtung}

Für das Entlüften selbst wurde eine Vorrichtung entwickelt, die schnelles und
blasenfreies Entlüften ermöglicht. Dazu wurde ein leerer Wasserbehälter
aufgeschnitten und mit zwei Schlauchanschlüssen versehen. An diesen sind die
Verbindungsschläuche zum Kühlkreis befestigt. Für ein blasenfreies Entlüften
müssen beide Anschlüsse stets unter Wasser stehen, und die Schläuche dürfen
keine Luft enthalten.

% FIGURE-VORSCHLAG 18: Foto der Entlüftungsvorrichtung (aufgeschnittener Behälter
% mit beiden Schlauchanschlüssen), ggf. mit Pfeilen für die Strömungsrichtung.

% TODO (optional): Nummerierter Ablauf des Entlüftens als \begin{enumerate}
% (Vorbefüllen -> Vorrichtung anschließen -> Pumpe starten -> Luftfreiheit prüfen).